% \iffalse meta-comment
% A LaTeX Class for Jacobs University Slides
% Copyright (c) 2007 Michael Kohlhase, all rights reserved
%               this file is released under the
%               Gnu Library Public Licences (LGPL)
%
% The original of this file is in the public repository at 
% http://github.com/KWARC/sTeX/
% \fi
% 
% \iffalse
%<cls>\NeedsTeXFormat{LaTeX2e}[1999/12/01]
%<cls>\ProvidesClass{mikoslides}[2014/01/07 v1.0 Jacobs slides]
%
%<*driver>
\documentclass{ltxdoc}
\usepackage{url,array,float,xspace}
\usepackage[show]{ed}
\usepackage{graphicx}
\usepackage[hyperref=auto,style=alphabetic]{biblatex}
\addbibresource{kwarcpubs.bib}
\addbibresource{extpubs.bib}
\addbibresource{kwarccrossrefs.bib}
\addbibresource{extcrossrefs.bib}
\makeindex
\floatstyle{boxed}
\newfloat{exfig}{thp}{lop}
\floatname{exfig}{Example}
\def\tracissue#1{\cite{sTeX:online},
\hyperlink{http://trac.kwarc.info/sTeX/ticket/#1}{issue #1}}
\usepackage{hyperref}
\begin{document}\DocInput{mikoslides.dtx}\end{document}
%</driver>
% \fi
% 
% \iffalse\CheckSum{380}\fi
% 
% \changes{v0.1}{2005/12/06}{Initial Version}
% \changes{v0.2}{2006/01/11}{course notes back on seminar}
% \changes{v0.3}{2007/05/16}{changing to Jacobs logo}
% \changes{v0.4}{2007/10/16}{re-basing the whole thing on beamer}
% \changes{v0.4}{2008/09/06}{moving line-end-comment to \texttt{omdoc.dtx}}
% \changes{v0.5}{2009/06/17}{eliminating mytwocolumns, this is better done by \texttt{beamer.cls}}
% \changes{v0.9}{2010/06/18}{this is almost done}
% \changes{v0.9}{2012/09/17}{basic options handling for the \texttt{frame} environment in
% notes mode}
% \changes{v0.9}{2013/08/23}{numbered sectocframes}
% \changes{v1.0}{2014/01/07}{adding \texttt{\textbackslash frameimage}}
% 
% \GetFileInfo{mikoslides.cls}
% \MakeShortVerb{\|}
%
% \def\twin#1#2{\index{#1!#2}\index{#2!#1}}
% \def\twintoo#1#2{{#1 #2}\twin{#1}{#2}}
% \def\atwin#1#2#3{\index{#1!#2!#3}\index{#3!#2 (#1)}}
% \def\atwintoo#1#2#3{{#1 #2 #3}\atwin{#1}{#2}{#3}}
%
% \def\scsys#1{{{\sc #1}}\index{#1@{\sc #1}}}
% \def\stex{\hbox{\raisebox{-.5ex}S\kern-.5ex\TeX}}
% \def\sTeX{\stex}
% \def\cnxml{\scshape{CNXml}}
% \def\connexions{\scshape{Connexions}}
% \def\element#1{{\ttfamily{#1}}}
% \def\snippet#1{{\ttfamily{#1}}}
% \def\cnxlatex{CNX\LaTeX\xspace}
% \def\mathml{{\scshape{MathML}}\xspace}
% \def\latexml{\hbox{{\LaTeX}ML}\xspace}
% \def\omdoc{OMDoc\xspace}
% \def\activemath{{\scshape{ActiveMath}}\xspace}
% \def\textwarning{\includegraphics[width=1.2em]{dangerous-bend}\xspace}
% 
% \title{Slides and Course Notes for Jacobs University\thanks{Version {\fileversion}
% (last revised {\filedate})}}
%    \author{Michael Kohlhase\\
%            Jacobs University, Bremen\\
%            \url{http://kwarc.info/kohlhase}}
% \maketitle
%
% \begin{abstract}
%   We present a document class from which we can generate both course slides and course
%   notes in a transparent way. Furthermore, we present a set of {\latexml} bindings for
%   these, so that we can also generate {\omdoc}-based course materials, e.g. for
%   inclusion in the {\activemath} system.
% \end{abstract}
%
% \tableofcontents\newpage
%
%\section{Introduction}
%
% This Document class is derived from |beamer.cls|~\cite{beamerclass:on}, specializes it
% with Jacobs stuff and adds a notes version that is more suited to printing than the one
% supplied by |beamer.cls|.
%
% 
%\section{The User Interface}\label{sec:user}
%
% The |mikoslides| class takes the notion of a slide frame from Till Tantau's excellent
% |beamer| class and adapts its notion of frames for use in the \sTeX and \omdoc. To
% support semantic course notes, it extends the notion of mixing frames and explanatory
% text, but rather than treating the frames as images (or integrating their contents into
% the flowing text), the |mikoslides| package displays the slides as such in the course
% notes to give students a visual anchor into the slide presentation in the course (and to
% distinguish the different writing styles in slides and course notes).
% 
% In practice we want to generate two documents from the same source: the slides for
% presentation in the lecture and the course notes as a narrative document for home
% study. To achieve this, the |mikoslides| class has two modes: \emph{slides mode} and
% \emph{notes mode} which are determined by the package option. 
%
% \subsection{Package Options}\label{sec:user:options}
% 
% The |mikoslides| class takes a variety of class options:\ednote{leaving out noproblems
% for the moment until we decide what to do with it.}
% \begin{compactitem}
% \item The options \DescribeMacro{slides}|slides|\DescribeMacro and {notes}|notes| switch
%   between slides mode and notes mode (see Section~\ref{sec:user:notesslides}).
% \item If the option \DescribeMacro{sectocframes}|sectocframes| is given, then special
%   frames with section table of contents are produced headers \ednote{document the functionality}
% \item \DescribeMacro{showmeta}|showmeta|. If this is set, then the metadata keys are
%   shown (see~\cite{Kohlhase:metakeys:ctan} for details and customization options).
% \item If the option \DescribeMacro{frameimages}|frameimages| is set, then slide mode
%   also shows the |\frameimage|-generated frames. 
% \end{compactitem}
% 
% \subsection{Notes and Slides}\label{sec:user:notesslides}
% 
% Slides are represented with the \DescribeEnv{frame}|frame| just like in the |beamer|
% class, see~\cite{Tantau:ugbc} for details. The |mikoslides| class adds the
% \DescribeEnv{note}|note| environment for encapsulating the course note
% fragments.\footnote{MK: it would be very nice, if we did not need this environment, and
% this should be possible in principle, but not without intensive LaTeX trickery. Hints to
% the author are welcome.} 
%
% \textwarning Note that it is essential to start and end the |notes| environment at the
% start of the line -- in particular, there may not be leading blanks -- else {\LaTeX}
% becomes confused and throws error messages that are difficult to decipher.
%
% \begin{exfig}
% \begin{verbatim}
% \begin{note}
%   We start this course with ...
% \end{note}
%
% \begin{frame}
%   \frametitle{The first slide}
%   ...
% \end{frame}
% \begin{note}
%   ... and more explanatory text
% \end{note}
%
% \begin{frame}
%   \frametitle{The second slide}
%   ...
% \end{frame}
% ...
% \end{verbatim}
% \caption{A typical Course Notes File}\label{fig:notesfile}
% \end{exfig}
% 
% By interleaving the |frame| and |note| environments, we can build course notes as shown
% in Figure~\ref{fig:notesfile}. 
% 
% Sometimes, we want to integrate slides as images after all -- e.g. because we already
% have a PowerPoint presentation, to which we want to add \sTeX notes. In this case we can
% use \DescribeMacro{\frameimage}|\frameimage[|\meta{opt}|]{|\meta{path}|}|, where
% \meta{opt} are the options of |\includegraphics| from the |graphicx|
% package~\cite{CarRah:tpp99} and \meta{path} is the file path (extension can be left off
% like in |\includegraphics|).
% 
% \subsection{Header and Footer Lines}\label{sec:user:headfootlines}
%
% \subsection{Colors and Highlighting}\label{sec:user:highlighting}
% The \DescribeMacro{\textwarning}|\textwarning| macro generates a warning
% sign: \textwarning
%
% \subsection{Front Matter, Titles, etc}\label{sec:user:matter}
%
% \subsection{Miscellaneous}\label{sec:user:misc}
%
% \subsection{Support for \textsf{MathHub}}\label{sec:user:mathhub}
% 
% Much of the \sTeX content is hosed on \textsf{MathHub} (\url{http://MathHub.info}), a
% portal and archive for flexiformal mathematics. \textsf{MathHub} offers GIT repositories
% (public and private escrow) for mathematical documentation projects, online and offline
% authoring and document development infrastructure, and a rich, interactive reading
% interface. The |modules| package supports repository-sensitive operations on
% \textsf{MathHub}.
% 
% Note that \textsf{MathHub} has two-level repository names of the form
% \meta{group}|/|\meta{repo}, where \meta{group} is a \textsf{MathHub}-unique repository
% group and \meta{repo} a repository name that is \meta{group}-unique. The file and
% directory structure of a repository is arbitrary -- except that it starts with the
% directory |source| because they are Math Archives in the sense
% of~\cite{HorIacJuc:cscpnrr11}. But this structure can be hidden from the \sTeX author
% with \textsf{MathHub}-enabled versions of the |modules| macros.
% 
% The \DescribeMacro{\mhframeimage}|\mhframeimage| macro is a variant of |\frameimage|
% with repository support. Instead of writing
% \begin{verbatim}
% \defpath{MathHub}{/user/foo/lmh/MathHub}
% \frameimage{\MathHub{fooMH/bar/source/baz/foobar}}
% \end{verbatim}
% we can simply write (assuming that |\MathHub| is defined as above)
% \begin{verbatim}
% \mhframeimage[fooMH/bar]{baz/foobar}
% \end{verbatim}
% Note that the |\mhframeimage| form is more semantic, which allows more advanced document
% management features in \textsf{MathHub}.
% 
% If |baz/foobar| is the ``current module'', i.e. if we are on the \textsf{MathHub} path
% \ldots|MathHub/fooMH/bar|\ldots, then stating the repository in the first optional
% argument is redundant, so we can just use
% \begin{verbatim}
% \mhframeimage{baz/foobar}
% \end{verbatim}
% Of course, neither {\LaTeX} nor \latexml know about the repositories when they are
% called from a file system, so we can use the |\mhcurrentrepos| macro from the |modules|
% package to tell them. But this is only needed to initialize the infrastructure in the
% driver file. In particular, we do not need to set it in in each module, since the
% |\importmhmodule| macro sets the current repository automatically.
% 
% \paragraph{Caveat} if you want to use the \textsf{MathHub} support macros (let's call
% them mh-variants), then every time a module is imported or a document fragment is
% included from another repos, the mh-variant |\importmhmodule| must be used, so that the
% ``current repository'' is set accordingly. To be exact, we only need to use mh-variants,
% if the imported module or included document fragment use mh-variants.
%
% \section{Limitations}\label{sec:limitations}
% 
% In this section we document known limitations. If you want to help alleviate them,
% please feel free to contact the package author. Some of them are currently discussed in
% the \sTeX TRAC~\cite{sTeX:online}. 
% \begin{compactenum}
% \item the class should be divided into concerns. \tracissue{1684}
% \item when option |book| or |report| is given together with |sectocframes| chapter-level
%   omgroups generate a spurious slide with a bare heading. This has something to do with
%   the fact that beamer does not support |\chapter|
% \end{compactenum}
% 
% \StopEventually{\newpage\PrintIndex\newpage\PrintChanges\printbibliography}
% 
%\section{The Implementation}\label{sec:impl}
%
% The |mikoslides| package generates two files: the {\LaTeX} package (all the code between
% {\textsf{$\langle$*package$\rangle$}} and {\textsf{$\langle$/package$\rangle$}}) and the
% {\latexml} bindings (between {\textsf{$\langle$*ltxml$\rangle$ and
%     $\langle$/ltxml$\rangle$}}). We keep the corresponding code fragments together,
% since the documentation applies to both of them and to prevent them from getting out of
% sync.
%
%\subsection{Initialization and Class Options}\label{sec:impl:init}
%
% For the {\latexml} bindings, we make sure the right perl packages are loaded.
%
%    \begin{macrocode}
%<*ltxml>
# -*- CPERL -*-
package LaTeXML::Package::Pool;
use strict;
use LaTeXML::Package;
DeclareOption('showmeta', sub {PassOptions('metakeys','sty',ToString(Digest(T_CS('\CurrentOption')))); });
DeclareOption('defindex', sub {PassOptions('statements','sty',ToString(Digest(T_CS('\CurrentOption')))); });
DeclareOption('notes', '');
DeclareOption('slides', '');
DeclareOption('noproblems', '');
DeclareOption('sectocframes', '');
DeclareOption('frameimages', '');
DeclareOption('report', sub {PassOptions('omdoc','cls',ToString(Digest(T_CS('\CurrentOption')))); });
DeclareOption('book', sub {PassOptions('omdoc','cls',ToString(Digest(T_CS('\CurrentOption')))); });
DeclareOption(undef, sub {PassOptions('omdoc','cls',ToString(Digest(T_CS('\CurrentOption'))));                          PassOptions('beamer','cls',ToString(Digest(T_CS('\CurrentOption'))));});
ProcessOptions();
%</ltxml>
%    \end{macrocode}
%
% For {\LaTeX} we define some Package Options and switches for the |mikoslides| class and activate
% them by passing them on to |beamer.cls| the appropriate packages.
%
%    \begin{macrocode}
%<*cls>
\DeclareOption{showmeta}{\PassOptionsToPackage{\CurrentOption}{metakeys}}
\DeclareOption{defindex}{\PassOptionsToPackage{\CurrentOption}{statements}}
\newif\ifnotes\notesfalse
\newif\ifsectocframes\sectocframesfalse 
\newif\ifframeimages\frameimagesfalse 
\newif\ifproblems\problemstrue
\DeclareOption{notes}{\notestrue}
\DeclareOption{slides}{\notesfalse}
\DeclareOption{noproblems}{\problemsfalse}
\DeclareOption{sectocframes}{\sectocframestrue}
\DeclareOption{frameimages}{\frameimagestrue}
%    \end{macrocode}
% the next two define the |frontmatter| environment so that the later |\renewcommand| does
% not lead to trouble.
%    \begin{macrocode}
\newif\if@part\@partfalse
\DeclareOption{report}{\@parttrue\PassOptionsToClass{\CurrentOption}{omdoc}}
\DeclareOption{book}{\@parttrue\PassOptionsToClass{\CurrentOption}{omdoc}}
\DeclareOption*{\PassOptionsToClass{\CurrentOption}{omdoc}
                \PassOptionsToClass{\CurrentOption}{beamer}}
\ProcessOptions
%</cls>
%<*ltxml>
RawTeX('\newif\ifnotes\notesfalse');
RawTeX('\newif\ifproblems\problemsfalse');
%</ltxml>
%    \end{macrocode}
%
% Depending on the options, we either load the |article|-based |omdoc| or the |beamer| class. In the first
% case, we also have to make the |beamer|-specific things available to |article| via the
% |beamerarticle| package. We use options to avoid loading theorem-like environments,
% since we want to use our own from the $\sTeX$ packages. 
%    \begin{macrocode}
%<*cls>
\ifnotes
  \LoadClass{omdoc}
  \RequirePackage{a4wide}
  \RequirePackage{marginnote}
  \RequirePackage{mdframed}
  \RequirePackage[notheorems,noamsthm,noxcolor]{beamerarticle}
\else
  \LoadClass[notheorems,noamsthm,10pt]{beamer}
  \newcounter{Item}
  \newcounter{paragraph}
  \newcounter{subparagraph}
  \newcounter{Hfootnote}
  \usetheme{Jacobs}
\fi
%</cls>
%<*ltxml>
LoadClass('omdoc');
RequirePackage('tikzinput');
DefConstructor('\usetheme{}','');
%</ltxml>
%    \end{macrocode}
% 
% now, we load the remaining packages for both versions. \ednote{MK: eventually (when tikz
% support is fully realized in \latexml) get rid of the standalone package} 
%
%    \begin{macrocode}
%<*cls>
\RequirePackage{stex}
\RequirePackage{tikzinput}
\RequirePackage{latexml}
\RequirePackage{amssymb}
\usepgflibrary{shapes}
\usetikzlibrary{arrows}
\usetikzlibrary{positioning}
\usetikzlibrary{tikzmark}%experimental/beta but very useful
\usetikzlibrary{fit}
\RequirePackage{url}
\RequirePackage{amsmath}
\RequirePackage{comment}
\RequirePackage{standalone}
\RequirePackage{textcomp}
%</cls>
%<*ltxml>
RequirePackage('stex'); 
RequirePackage('latexml');
RequirePackage('amssymb');
RequirePackage('graphicx');
RequirePackage('tikz');
RequirePackage('url');
 RequirePackage('amsmath');
%</ltxml>
%    \end{macrocode}
% 
% \subsection{Notes and Slides}\label{sec:impl:noteslides}
%
% We define the sizes of slides in the notes. Somehow, we cannot get by with the same
% here. 
%
%    \begin{macrocode}
%<*cls>
\newcounter{slide}
\newlength{\slidewidth}\setlength{\slidewidth}{12.8cm}
\newlength{\slideheight}\setlength{\slideheight}{9cm}
%</cls>
%<*ltxml>
DefRegister('\slidewidth'     => Dimension('13.6cm'));
DefRegister('\slideheight'     => Dimension('9cm'));
%</ltxml>
%    \end{macrocode}
% 
% \begin{environment}{note}
% The |note| environment is used to leave out text in the |slides| mode. It does not have
% a counterpart in OMDoc. So for course notes, we define the |note| environment to be a
% no-operation otherwise we declare the |note| environment as a comment via the |comment|
% package.
%    \begin{macrocode}
%<*cls>
\ifnotes%
  \renewenvironment{note}{%
    \ignorespaces%
  }{}%
\else%
  \excludecomment{note}%
\fi%
%</cls>
%<*ltxml>
DefEnvironment('{note}','#body');
%</ltxml>
%    \end{macrocode}
% \end{environment}
% 
%
% We start by giving the {\latexml} binding for the |frame| environment from the |beamer|
% class. We first set up the slide boxes in |article| mode. We set up sizes and provide a
% box register for the frames and a counter for the slides.
% 
%    \begin{macrocode}
%<*cls>
\ifnotes
  \newlength{\slideframewidth}
  \setlength{\slideframewidth}{1.5pt}
%    \end{macrocode}
% 
% \begin{environment}{frame}
%   We first define the keys. 
%    \begin{macrocode}
  \addmetakey{frame}{label}
  \addmetakey[yes]{frame}{allowframebreaks}
  \addmetakey{frame}{allowdisplaybreaks}
  \addmetakey[yes]{frame}{fragile}
  \addmetakey[yes]{frame}{shrink}
  \addmetakey[yes]{frame}{squeeze}
%    \end{macrocode}
%   We redefine the |itemize| environment so that it looks more like the one in |beamer|
%   with |Jacobs| theme. We create the box with the |mdframed| environment from the
%   equinymous package.
%   Then we define the environment, read them, and construct the slide number and label. 
%    \begin{macrocode}
  \renewenvironment{frame}[1][]{%
    \metasetkeys{frame}{#1}%
    \stepcounter{slide}%
    \def\@currentlabel{\theslide}%
    \ifx\frame@label\@empty%
    \else%
      \label{\frame@label}%
    \fi%
%    \end{macrocode}
%   We redefine the |itemize| environment so that it looks more like the one in |beamer|
%   with |Jacobs| theme. 
%    \begin{macrocode}
    \def\itemize@level{outer}%
    \def\itemize@outer{outer}%
    \def\itemize@inner{inner}%
    \renewcommand\newpage{}%
    \renewcommand\metakeys@show@keys[2]{\marginnote{{\scriptsize ##2}}}%
    \renewenvironment{itemize}{%
      \ifx\itemize@level\itemize@outer%
        \def\itemize@label{$\rhd$}%
      \fi%
      \ifx\itemize@level\itemize@inner%
        \def\itemize@label{$\scriptstyle\rhd$}%
      \fi%
      \begin{list}%
      {\itemize@label}%
      {\setlength{\labelsep}{.3em}%
       \setlength{\labelwidth}{.5em}%
       \setlength{\leftmargin}{1.5em}%
      }%
      \edef\itemize@level{\itemize@inner}%
    }{%
      \end{list}%
    }%
%    \end{macrocode}
% We create the box with the |mdframed| environment from the equinymous package.
%    \begin{macrocode}
    \begin{mdframed}[linewidth=\slideframewidth,skipabove=1ex,skipbelow=1ex,userdefinedwidth=\slidewidth,align=center]\sf%
  }{%
    \medskip\miko@slidelabel\end{mdframed}%
  }%
%</cls>
%<*ltxml>
DefEnvironment('{frame}[]',
   "<omdoc:omgroup layout='slide'>"
  .  "#body\n"
  ."</omdoc:omgroup>\n\n",
 afterDigestBegin=>sub {
   $_[1]->setProperty(theory=>LookupValue('current_module')); });
%</ltxml>#$
%    \end{macrocode}
% \end{environment}
% 
% Now, we need to redefine the frametitle (we are still in course notes mode). 
% \begin{macro}{\frametitle}
%    \begin{macrocode}
%<*cls>
  \renewcommand{\frametitle}[1]{{\Large\bf\sf\color{blue}{#1}}\medskip}%
\fi %ifnotes
%</cls>
%<*ltxml>
DefConstructor('\frametitle{}',
  "\n<omdoc:metadata><dc:title>#1</dc:title></omdoc:metadata>");
%</ltxml>
%    \end{macrocode}
% \end{macro}
% 
% \begin{macro}{\frameimage}
%   We have to make sure that the width is overwritten, for that we check the
%   |\Gin@ewidth| macro from the |graphicx| package\ednote{MK@DG; we need to do that in
%   the LaTeXML binding as well!}
%    \begin{macrocode}
%<*cls>
\newrobustcmd\frameimage[2][]{%
  \stepcounter{slide}%
  \ifframeimages%
    \def\Gin@ewidth{}\setkeys{Gin}{#1}%
    \ifnotes%
    \else%
      \vfill%
    \fi%
    \ifx\Gin@ewidth\@empty%
      \mycgraphics[width=\slidewidth,#1]{#2}\else\mycgraphics[#1]{#2}%
    \fi%
    \par\strut\hfill{\footnotesize Slide \arabic{slide}}%
    \ifnotes%
    \else%
      \vfill%
    \fi%
  \fi%
}% ifframeimages
%</cls>
%<*ltxml>
DefMacro('\frameimage[]{}','\@frameimage{\includegrahics[#1,width=\slidewidth]{#2}}');
DefConstructor('\@frameimage{}',"<omdoc:omgroup layout='slide'>#1</omdoc:omgroup>\n");
%</ltxml>
%    \end{macrocode}
% \end{macro}

% 
% \subsection{Header and Footer Lines}\label{sec:impl:headfootlines}
%
% Now, we set up the infrastructure for the footer line of the slides, we use boxes for
% the logos, so that they are only loaded once, that considerably speeds up processing.
% \begin{macro}{\setslidelogo}
% The default logo is the logo of Jacobs University. Customization can be done by |\setslidelogo{|\meta{logo name}|}|.
%    \begin{macrocode}
%<*cls>
\newlength{\slidelogoheight}
\ifnotes%
  \setlength{\slidelogoheight}{.4cm}%
\else%
  \setlength{\slidelogoheight}{1cm}%
\fi%
\newsavebox{\slidelogo}%
\sbox{\slidelogo}{\includegraphics[height=\slidelogoheight]{jacobs-logo}}%
\newrobustcmd{\setslidelogo}[1]{%
  \sbox{\slidelogo}{\includegraphics[height=\slidelogoheight]{#1}}%
}%
%    \end{macrocode}
% \end{macro}
%
% \begin{macro}{\setsource}
% |\source| stores the writer's name. By default it is {\it Michael Kohlhase} since he is the main user and designer of this package. |\setsource{|\meta{name}|}| can change the writer's name.
%    \begin{macrocode}
\def\source{Michael Kohlhase}% customize locally
\newrobustcmd{\setsource}[1]{\def\source{#1}}%
%    \end{macrocode}
% \end{macro}
%
% \begin{macro}{\setlicensing}
% Now, we set up the copyright and licensing. By default we use the Creative Commons Attribuition-ShareAlike license to strengthen the public domain. If package |hyperref| is loaded, then we can attach a hyperlink to the license logo. |\setlicensing[|\meta{url}|]{|\meta{logo name}|}| is used for customization, where ||\meta{url}|| is optional.
%    \begin{macrocode}
\def\copyrightnotice{\footnotesize\copyright:\hspace{.3ex}{\source}}%
\newsavebox{\cclogo}%
\sbox{\cclogo}{\includegraphics[height=\slidelogoheight]{cc_somerights}}%
\newif\ifcchref\cchreffalse%
\AtBeginDocument{%
  \@ifpackageloaded{hyperref}{\cchreftrue}{\cchreffalse}
}%
\def\licensing{%
  \ifcchref%
    \href{http://creativecommons.org/licenses/by-sa/2.5/}{\usebox{\cclogo}}%
  \else%
    {\usebox{\cclogo}}%
  \fi%
}%
\newrobustcmd{\setlicensing}[2][]{%
  \def\@url{#1}%
  \sbox{\cclogo}{\includegraphics[height=\slidelogoheight]{#2}}%
  \ifx\@url\@empty%
    \def\licensing{{\usebox{\cclogo}}}%
  \else%
    \def\licensing{%
	  \ifcchref%
	    \href{#1}{\usebox{\cclogo}}%
	  \else%
	    {\usebox{\cclogo}}%
	  \fi%
    }%
  \fi%
}%
%    \end{macrocode}
% \end{macro} 
%
% \begin{macro}{\slidelabel}
% Now, we set up the slide label for the |article| mode.\ednote{see that we can use the themes for the slides some day. This is all fake.}
%    \begin{macrocode}
\newrobustcmd\miko@slidelabel{%
  \vbox to \slidelogoheight{%
    \vss\hbox to \slidewidth%
    {\licensing\hfill\copyrightnotice\hfill\arabic{slide}\hfill\usebox{\slidelogo}}%
  }%
}%
%    \end{macrocode}
% \end{macro}
% 
% \subsection{Colors and Highlighting}\label{sec:impl:highlighting}
%
% Now, we set up an infrastructure for highlighting phrases in slides. Note that we use
% content-oriented macros for highlighting rather than directly using color markup. 
% The first thing to to is to adapt the green so that it is dark enough for most beamers
%    \begin{macrocode}
\AtBeginDocument{%
  \definecolor{green}{rgb}{0,.5,0}%
  \definecolor{purple}{cmyk}{.3,1,0,.17}%
}%
%    \end{macrocode}
%
% We customize the |\defemph|, |\notemph|, and |\stDMemph| macros with colors for the use
% in the |statements| package. Furthermore we customize the |\@@lec| macro for the
% appearance of line end comments in |\lec|.
%
%    \begin{macrocode}
% \def\STpresent#1{\textcolor{blue}{#1}}
\def\defemph#1{{\textcolor{magenta}{#1}}}
\def\notemph#1{{\textcolor{magenta}{#1}}}
\def\stDMemph#1{{\textcolor{blue}{#1}}}
\def\@@lec#1{(\textcolor{green}{#1})}
%</cls>
%<*ltxml>
#DefMacro('\defemph{}','{\textcolor{magenta}{#1}}');
#DefMacro('\notemph{}','{\textcolor{magenta}{#1}}');
%</ltxml>
%    \end{macrocode}
%
% I like to use the dangerous bend symbol for warnings, so we provide it here.
% \begin{macro}{\textwarning}
%   as the macro can be used quite often we put it into a box register, so that it is only
%   loaded once. 
%    \begin{macrocode}
%<*cls>
\pgfdeclareimage[width=.9em]{miko@small@dbend}{dangerous-bend}
\def\smalltextwarning{%
  \pgfuseimage{miko@small@dbend}%
  \xspace%
}%
\pgfdeclareimage[width=1.5em]{miko@dbend}{dangerous-bend}
\newrobustcmd\textwarning{%
  \raisebox{-.05cm}{\pgfuseimage{miko@dbend}}%
  \xspace%
}%
\pgfdeclareimage[width=2.5em]{miko@big@dbend}{dangerous-bend}%
\newrobustcmd\bigtextwarning{%
  \raisebox{-.05cm}{\pgfuseimage{miko@big@dbend}}%
  \xspace%
}%
%</cls>
%<*ltxml>
DefMacro('\textwarning','\@textwarning\xspace');
DefConstructor('\@textwarning',"");
%</ltxml>
%    \end{macrocode}
% \end{macro}
% 
% \subsection{Front Matter, Titles, etc}
%
% We need to redefine the frontmatter macros inherited from the |beamer| class for
% LaTeXML, since there they take an optional argument.
%    \begin{macrocode}
%<*ltxml>
DefMacro('\title[]{}', '\@add@frontmatter{ltx:title}{#1}');
DefMacro('\date[]{}', '\@add@frontmatter{ltx:date}[role=creation]{#1}');
DefMacro('\author[]{}', sub { andSplit(T_CS('\@author'),$_[1]); });#$
%</ltxml>
%    \end{macrocode}
%
% Now, we specialize the slide environment that we have implemented above or inherited
% from |seminar.cls| for some abbreviations, e.g. separator slides and title slides.
%    \begin{macrocode}
%<*cls>
\ifnotes%
  \newrobustcmd\titleframe{\maketitle}%
\else%
  \newrobustcmd\titleframe{%
    \begin{frame}%
    \titlepage%
    \end{frame}%
  }%
\fi%
\newenvironment{titleframewith}{%
  \begin{frame}%
  \titlepage%
}{%
  \end{frame}%
}%
\newenvironment{ttitle}{%
  \begin{center}%
  \LARGE%
  \begin{tabular}{|c|}%
  \hline%
}{%
  \\\hline%
  \end{tabular}%
  \end{center}%
  \vspace{1ex minus 1ex}%
}%
\newenvironment{ttitlejoint}[1]{%
  \newbox\boxwith%
  \setbox\boxwith\hbox{%
    \begin{tabular}{c}{\em joint work with}\\#1\end{tabular}%
  }%
  \begin{center}%
  \LARGE%
  \begin{tabular}{c}%
  \color{red}%
}{%
  \\\box\boxwith%
  \end{tabular}%
  \end{center}%
  \vspace{1ex minus 1ex}%
}%
%</cls>
%<*ltxml>
DefConstructor('\titleframe',"<omdoc:ignore>titleframe elided here</omdoc:ignore>");
DefEnvironment('{titleframewith}',
               "<omdoc:ignore>begin elided titleframe</omdoc:ignore>"
               . "#body"
              ."<omdoc:ignore>end elided titleframe</omdoc:ignore>"); 
DefEnvironment('{titleslide}',"");
DefEnvironment('{titleslide}',"<omdoc:omgroup>#body</omdoc:omgroup>");
DefEnvironment('{ttitle}', "\n<dc:title>#body</dc:title>");
%</ltxml>
%    \end{macrocode}
%
%    \begin{macrocode}
%      Must be first command on slide to make positioning work.
%<*cls>
\newrobustcmd\putgraphicsat[3]{%
  \begin{picture}(0,0)\put(#1){\includegraphics[#2]{#3}}\end{picture}%
}%
\newrobustcmd\putat[2]{%
  \begin{picture}(0,0)\put(#1){#2}\end{picture}%
}%
%</cls>
%    \end{macrocode}
%
% \subsection{Sectioning}
% If the |sectocframes| option is set, then we make section frames. 
%    \begin{macrocode}
%<*cls>
\ifsectocframes%
  \if@part%
    \newcounter{mpart}%
    \newcounter{mchapter}%
    \newcounter{msection}[mchapter]%
  \else%
    \newcounter{msection}%
  \fi%
  \newcounter{msubsection}[msection]%
  \newcounter{msubsubsection}[msubsection]%
  \newcounter{msubsubsubsection}[msubsubsection]%
  \ifnotes%
  \else% only in slides
    \renewcommand\at@begin@omgroup[3][]{%
      \begin{frame}%
      \vfill\Large\centering%
      \red{%
      	\ifcase\section@level\or%
        \stepcounter{mpart}Part \Roman{mpart}\or%
        \stepcounter{mchapter}Chapter \arabic{mchapter}\or
        \stepcounter{msection}\if@part\arabic{mchapter}.\fi\arabic{msection}\or
        \stepcounter{msubsection}\if@part\arabic{mchapter}.\fi\arabic{msection}.\arabic{msubsection}\or
        \stepcounter{msubsubsection}\if@part\arabic{mchapter}.\fi\arabic{msection}.\arabic{msubsection}.\arabic{msubsubsection}.\or
        \stepcounter{msubsubsubsection}\if@part\arabic{mchapter}.\fi\arabic{msection}.\arabic{msubsection}.\arabic{msubsubsection}.\arabic{msubsubsubsection}\fi%
        \quad #3%
      }%
      \vfill%
      \end{frame}%
    }%
  \fi% ifnotes
\fi% ifsectocframes
%</cls>
%    \end{macrocode}
%
% \subsection{Miscellaneous}
%
% The following fixes an error I do not understand, this has something to do with
% beamer compatibility, which has similar definitions but only up to 1. 
%    \begin{macrocode}
%<*cls>
\expandafter\def\csname Parent2\endcsname{}
%    \begin{macrocode}
%
% We need to disregard the columns macros introduced by the |beamer| class
%    \begin{macrocode}
\ifnotes%
  \renewenvironment{columns}{%
    \par\noindent%
    \begin{minipage}%
    \slidewidth\centering\leavevmode%
  }{%
    \end{minipage}\par\noindent%
  }%
  \newsavebox\columnbox%
  \renewenvironment{column}[1]{%
    \begin{lrbox}{\columnbox}\begin{minipage}{#1}%
  }{%
    \end{minipage}\end{lrbox}\usebox\columnbox%
  }%
\fi%
%</cls>
%<*ltxml>
DefEnvironment('{columns}',"#body");
DefEnvironment('{column}{}',"#body");
%    \end{macrocode}
%
% We also need to deal with overlay specifications introduced by the |beamer|
% class.\ednote{this is just to keep latexml quiet, no real functionality here.}
%
% \ednote{Deyan: We reuse the CMP itemizations defined in the omdoc.cls.ltxml binding,
%   adjusting the parameters to be overlay-sensitive }
%    \begin{macrocode}
DefConstructor('\uncover','#1');
#Define a Beamer Overlay Parameter type
 DefParameterType('BeamerOverlay', sub {
   my ($gullet) = @_;
   my $tok = $gullet->readXToken;
   if (ref $tok && ToString($tok) eq '<') {
     $gullet->readUntil(T_OTHER('>'));
   } else {
     $gullet->unread($tok) if ref $tok;
     undef; }},
     reversion=> sub {
 (T_OTHER('<'), $_[0]->revert, T_OTHER('>'));
       });

#Take the "from" field of the overlay range
sub overlayFrom {
  return "" unless defined $_[0];
  my $overlay=ToString($_[0]); $overlay =~ /^(\d+)/; $1;}

#Reuse the CMP itemizations, only adjust the \item constructors.
DefMacro('\beamer@group@item[] OptionalBeamerOverlay IfBeginFollows', sub {
  my($gullet,$tag,$overlay,$needwrapper)=@_;
  $overlay=$overlay||T_OTHER("");
  ( T_CS('\group@item@maybe@unwrap'),
    ($needwrapper ? (Invocation(T_CS('\beamer@group@item@wrap'),$tag,$overlay)->unlist) : ()) ); });
DefConstructor('\beamer@group@item@wrap {} OptionalBeamerOverlay',
       "<omdoc:omtext ?#2(overlay='&overlayFrom(#2)')()>"
       . "?#1(<dc:title>#1</dc:title>)()"
             . "<omdoc:CMP>",
       beforeDigest=>sub {
 Let('\group@item@maybe@unwrap','\group@item@unwrap');
 #$_[0]->bgroup;
return; },
       properties=>sub{ RefStepItemCounter(); });
#DefConstructor('\beamer@itemize@item[] OptionalBeamerOverlay',
#       "<omdoc:li ?#2(overlay='&overlayFrom(#2)')() >"
#     . "?#1(<dc:title>#1</dc:title>)()",
#       properties=>sub{ RefStepItemCounter(); });
DefConstructor('\beamer@enumerate@item[] OptionalBeamerOverlay',
       "<omdoc:li ?#2(overlay='&overlayFrom(#2)')() >"
     . "?#1(<dc:title>#1</dc:title>)()",
       properties=>sub{ RefStepItemCounter(); });
DefConstructor('\beamer@description@item[] OptionalBeamerOverlay',
       "<omdoc:di ?#2(overlay='&overlayFrom(#2)')() >"
       . "?#1(<omdoc:dt>#1</omdoc:dt>)()<omdoc:dd>", # trust di and dt to autoclose
       properties=>sub{ RefStepItemCounter(); });
%</ltxml>#$
%    \end{macrocode}
% Now, some things that are imported from the |pgf| and |beamer| packages: 
%    \begin{macrocode}
%<*ltxml>
DefMacro('\putgraphicsat{}{}{}','\mygraphics[#2]{#3}');
DefMacro('\putat{}{}','#2');
%</ltxml>
%    \end{macrocode}
% 
%
%    \begin{macrocode}
%<*cls>
\ifproblems%
  \newenvironment{problems}{}{}%
\else%
  \excludecomment{problems}%
\fi%
%</cls>
%<*ltxml>
DefEnvironment('{problems}','#body');
%</ltxml>
%    \end{macrocode}
%
% \subsection{Support for \textsf{MathHub}}\label{sec:user:mathhub}
% 
% \begin{macro}{\mhframeimage}
%   Use the current value of |\mh@currentrepos| or the value of the |mhrepos| key if it is
%   given in |\frameimage|.
%    \begin{macrocode}
%<cls>\addmetakey{Gin}{mhrepos}
%<ltxml>DefKeyVal('Gin','mhrepos','Semiverbatim');
%<ltxml>RawTeX('
%<*ltxml|cls>
\newrobustcmd\mhframeimage[2][]{%
  \metasetkeys{Gin}{#1}%
  \edef\mh@@repos{\mh@currentrepos}%
  \ifx\Gin@mhrepos\@empty%
    \frameimage[#1]{\MathHub{\mh@currentrepos/source/#2}}%
  \else%
    \frameimage[#1]{\MathHub{\Gin@mhrepos/source/#2}}%
  \fi%
}%
%</ltxml|cls>
%<ltxml>');
%    \end{macrocode}
% \end{macro}
%
% \subsection{Finale}
%
% Finally, we set the slide body font to the sans serif, and we terminate the {\latexml}
% bindings file with a success mark for perl.
%    \begin{macrocode}
%<cls>\ifnotes\else\sf\fi
%<ltxml>1;
%    \end{macrocode}
% \Finale
\endinput
% \iffalse
%%% Local Variables: 
%%% mode: doctex
%%% TeX-master: t
%%% End: 
% \fi
% LocalWords:  mikoslides dtx beamer omdoc notheorems noamsthm beamerarticle sc Licences
% LocalWords:  graphicx slidelabel stex latexml amssymb tikz url CPERL  amsmath filedate
% LocalWords:  LoadClass RequirePackage DefRegister DefEnvironment omgroup rgb frameimage
% LocalWords:  afterDigestBegin setProperty LookupValue DefConstructor hyperref Tantau's
% LocalWords:  cmyk lec DefMacro titleslide ttitle RawTeX metadata etoolbox cls emph ugbc
% LocalWords:  noproblems linkcolor bookmarksopen citecolor urlcolor colorlinks Tantau di
% LocalWords:  breaklinks plainpages pdfpagelabels srcref iffalse texttt atwin exfig dt
% LocalWords:  mytwocolumns twintoo atwintoo scsys sc hbox raisebox cnxml impl notesfile
% LocalWords:  scshape ttfamily cnxlatex mathml activemath fileversion newpage tpp99 dd
% LocalWords:  maketitle tableofcontents ednote compactitem showmeta showmeta beamerclass
% LocalWords:  sectocframes sectocframes textwarning textwarning compactenum eq omgroups
% LocalWords:  includegraphics tracissue printbibliography textsf langle textsf rangle
% LocalWords:  langle ltxml metakeys newif ifnotes notesfalse ifsectocframes rangle putat
% LocalWords:  sectocframesfalse ifproblems problemstrue notestrue marginnote frontmatter
% LocalWords:  problemsfalse sectocframestrue mdframed noxcolor newcounter ifx equinymous
% LocalWords:  Hfootnote usetheme tikzinput usepgflibrary usetikzlibrary rhd ignorespaces
% LocalWords:  tikzmark textcomp newlength slidewidth setlength slidewidth miko setkeys
% LocalWords:  slideheight slideheight renewenvironment excludecomment itenize Gin@ewidth
% LocalWords:  slideframewidth slideframewidth surroundwithmdframed addmetakey Gen@ewidth
% LocalWords:  allowframebreaks allowdisplaybreaks metasetkeys stepcounter sbox mpart baz
% LocalWords:  currentlabel theslide renewcommand scriptsize scriptstyle hspace mchapter
% LocalWords:  medskip linewidth skipabove skipbelow frametitle newenvironment msection
% LocalWords:  slidelogoheight newsavebox slidelogo slidelogo jacobs-logo vbox ifcase
% LocalWords:  Attribuition-ShareAlike copyrightnotice footnotesize cclogo vss minipage
% LocalWords:  cclogo somerights ifcchref cchreffalse ifpackageloaded usebox mycgraphics
% LocalWords:  cchreftrue usebox newcommand hfill hfill definecolor definecolor endinput
% LocalWords:  defemph notemph stDMemph STpresent textcolor textwarigrening ltx HorIacJuc
% LocalWords:  pgfdeclareimage dbend pgfuseimage xspace titleframe titlepage mycgraphics
% LocalWords:  titleframewith hline vspace ttitlejoint newbox boxwith boxwith msubsection
% LocalWords:  putgraphicsat beginomgroup ifnum vfill vfill noindent leavevmode cscpnrr11
% LocalWords:  columnbox lrbox needwrapper unlist omtext bgroup autoclose pgf readXToken
% LocalWords:  mygraphics doctex NeedsTeXFormat textbackslash userdefinedwidth foobar
%  LocalWords:  includegrahics smalltextwarning msubsubsection msubsubsubsection ldots
%  LocalWords:  mhframeimage mhframeimage mhcurrentrepos importmhmodule mh-variants Deyan
%  LocalWords:  mh-variant bigtextwarning setbox mh@currentrepos mhrepos Semiverbatim
%  LocalWords:  frameimages frameimages ifframeimages frameimagesfalse frameimagestrue
%  LocalWords:  expandafter csname endcsname specializes customization flexiformal colors
%  LocalWords:  initialize Initialization defindex realized itemize newrobustcmd
%  LocalWords:  specialize centering itemizations
